%% kapa: exam=1 round=36 year=2025 period=1 form=A subject=civil_law subject_ko=민법 seq=13 number=13
%% source: Q-Net 감정평가사 기출문제 (cst00309) artlSeq=5231525 file=raw/1008/5231525/2025년 제36회 감정평가사 1차시험 1교시 문제지 원본.pdf
%% conversion: kordoc 4.24.51 -> tools/gen1.py (auto); review=proofread-sonnet; flags=none
\begin{question}{13}
복대리에 관한 설명으로 옳지 않은 것은? (다툼이 있으면 판례에 따름)
\begin{choices}
  \item 임의대리인은 본인의 승낙이 있거나 부득이한 사유있는 때가 아니면 복대리인을 선임하지 못한다.
  \item 복대리인이 적법하게 선임되면 대리인의 대리권은 소멸한다.
  \item 대리인의 대리권이 소멸하면 복대리인의 대리권도 소멸한다.
  \item 복대리인은 제3자에 대하여 대리인과 동일한 권리의무가 있다.
  \item 복임권 없는 대리인에 의하여 선임된 복대리인의 권한도 권한을 넘은 표현대리의 기본대리권이 될 수 있다.
\end{choices}
\end{question}
\answer{2}
